\begin{minipage}[t]{0.74\linewidth}
En el quadrat de la figura, de 2,00~m de costat, hi ha dues càrregues $Q_1 = 9,00\ \mu\mathrm{C}$ i $Q_2 = -9,00\ \mu\mathrm{C}$ en els vèrtexs de l'esquerra.
\end{minipage}\hfill
\begin{minipage}[t]{0.2\linewidth}
\vspace{0pt}
\figura[1]{fig1.png}
\end{minipage}

\begin{apartat}{1}
Determineu la intensitat del camp elèctric en el centre del quadrat.
\end{apartat}

\begin{apartat}{1}
En el centre del quadrat hi situem una tercera càrrega $Q_3 = 7,00\ \mu\mathrm{C}$. Calculeu el treball que farà la força elèctrica que actua sobre $Q_3$ quan la traslladem del centre del quadrat al vèrtex inferior dret.
\end{apartat}

\blocetiquetat{Dada:}{$k = 9,00\times 10^{9}\ \mathrm{N\,m^2\,C^{-2}}$}
